\chapter{FX Market Making}\label{ch:hf:fx-market-making}

A non-bank liquidity provider streams euro-dollar prices to a hundred banks, brokers and platforms at once, each with its own price; a trade on one
of them is hedged, or not, within microseconds on another. In this chapter's model, the stream's last look lifts the capture on the most toxic
tier from 0.13 to 0.20 basis point a request and lowers it on the others, for a third of all requests rejected; and internalising the flow instead
of hedging it doubles the day's result, but only if the stream's skew brings offsetting trades back within a few hundred trades.

\section{Principal electronic liquidity provision}

Foreign exchange (One Quant Book 2, chapters 14 and 15) trades over the counter, in a web of venues: primary interdealer venues, electronic
communication networks, single-dealer and multi-dealer platforms, and the liquidity aggregators of One Quant Book 10, chapter 21. \omterm{def:hf:the-business-of-market-making:emm}{Electronic market
makers}, banks and non-banks, quote as principals: they stream prices, take the other side of every trade they accept, and manage the inventory
that results.

\begin{dated}{2026-09}{The size of the market}\label{dat:hf:fx-market-making:bis}
The BIS Triennial Central Bank Survey's preliminary results, published on 30 September 2025, put FX trading at \$9.6 trillion a day in April 2025.
FX swaps remained the most traded instrument at \$4 trillion a day; turnover in spot rose 42\% and in outright forwards 60\% since April 2022, to 31\%
and 19\% of the total. The US dollar was on one side of 89\% of all trades; sales desks in the United Kingdom, the United States, Singapore and Hong
Kong accounted for 75\%.
\end{dated}

Chaboud, Chiquoine, Hjalmarsson and Vega studied the rise of algorithmic trading in the interdealer FX market with data that identified
computer-generated activity. Algorithmic trading improved two measures of price efficiency, the frequency of triangular arbitrage opportunities and
the autocorrelation of high-frequency returns; fewer arbitrage opportunities came mainly from computers taking liquidity, which may impose higher
adverse selection costs on slower traders, and less autocorrelation from computers providing it.

\section{Pricing from many venues}

The provider's \omterm{def:hf:fair-value:fair}{fair price} (chapter 2) combines the venues it sees, each with its own noise and its own delay. Weighting each venue's mid by the
inverse of its variance, the noise plus the variance the rate has accumulated since that venue last updated
(\cref{lst:hf:fx-market-making:fair}), beats the primary venue alone: in the chapter's model of three venues (noise of 0.3, 0.6 and 0.8 basis point,
updating every 1, 3 and 5 milliseconds on average), the consolidated price is 0.26 basis point from the true rate, the primary 0.30.

\begin{definition}[Synthetic cross]\label{def:hf:fx-market-making:cross}
A \emph{synthetic cross}\index{synthetic cross} is a price for a currency pair built from two pairs that share a currency, such as euro-yen from
euro-dollar and dollar-yen, with its bid from the legs' prices that the provider can hit and its offer from those it can lift.
\end{definition}

With euro-dollar at 1.0850/1.0851 and dollar-yen at 150.20/150.21, the synthetic euro-yen is 162.967/162.993: 2.6 pips wide, the two legs' spreads
compounded. A provider that quotes the cross tighter than its synthetic must expect to internalise part of the flow, or hedge at the direct cross's
own market when that is cheaper.

\section{Streams, tiers and last look}

The provider streams a price to each client or platform, wider for tiers whose flow is more informed (One Quant Book 9, chapter 24 on client
tiering). Last look (One Quant Book 2, chapter 15, where the FX Global Code's principles for it are given) lets it hold a request for a few
milliseconds and reject it if the price has moved. The chapter's streams: tier 1 (banks and platforms, with fast arbitrage clients) at 0.3 basis
point, 30\% of its requests informed; tier 2 (brokers) at 0.5, 10\% informed; tier 3 (corporates) at 1.0, none. An informed request follows a
0.6-basis-point move the stream has not yet shown. Hold 5 milliseconds; reject asymmetrically any request whose price has moved more than 0.1 basis
point in the client's favour, or symmetrically in either direction (\cref{lst:hf:fx-market-making:stream}).

\begin{omfigure}
\begin{tikzpicture}[font=\footnotesize,>=stealth,box/.style={draw,rounded corners,fill=omThm!8,minimum height=0.7cm,align=center}]
\draw[->,thick] (0,0) -- (11.2,0) node[right] {time};
\foreach \x/\l in {0.6/{stream shows\\$S$},3.2/{request at $S$\\arrives},7.6/{check after\\5 ms}} \draw (\x,0.12) -- (\x,-0.12) node[below,align=center] {\l};
\draw[decorate,decoration={brace,amplitude=5pt}] (3.2,0.35) -- node[above=4pt] {hold window} (7.6,0.35);
\draw[omDef,thick] (0.6,1.3) -- (2.4,1.3) -- (2.4,1.9) -- (7.6,1.9);
\node[font=\scriptsize,anchor=west,omDef] at (2.45,1.6) {true rate moves 0.6 bp};
\node[font=\scriptsize,anchor=east] at (0.55,1.3) {rate};
\node[box] (acc) at (10.1,1.9) {accept: fill\\at $S$};
\node[box,fill=omDef!15] (rej) at (10.1,0.95) {reject: price moved\\$>0.1$ bp};
\draw[->] (7.7,1.3) -- (acc.west); \draw[->] (7.7,1.3) -- (rej.west);
\node[font=\scriptsize,align=left,anchor=north west] at (0.2,-1.05) {asymmetric: reject only moves in the client's favour (the informed request's);\\
symmetric: reject moves in either direction, the provider's favour too.};
\end{tikzpicture}

\omcaption{Last look as the chapter's stream applies it: the provider shows a price $S$; a client's request at $S$ arrives, informed ones after a
0.6-basis-point move of the true rate the stream has not yet shown; the provider holds it for 5 milliseconds and rejects it if the price has moved
by more than 0.1 basis point, only in the client's favour (asymmetric) or in either direction (symmetric). Source:
\texttt{firm.fxlp} on \texttt{firm.lastlook}.}\label{fig:hf:fx-market-making:lastlook}
\end{omfigure}

\begin{center}
\small
\begin{tabular}{@{}lrrrrrr@{}}
\toprule
& \multicolumn{2}{c}{no last look} & \multicolumn{2}{c}{asymmetric} & \multicolumn{2}{c}{symmetric}\\
\cmidrule(lr){2-3}\cmidrule(lr){4-5}\cmidrule(l){6-7}
tier & capture & rejects & capture & rejects & capture & rejects\\
\midrule
tier 1 (30\% informed) & 0.125 & 0 & 0.200 & 43.5\% & 0.140 & 56.8\%\\
tier 2 (10\% informed) & 0.453 & 0 & 0.395 & 26.5\% & 0.280 & 43.1\%\\
tier 3 (none) & 1.008 & 0 & 0.851 & 18.3\% & 0.637 & 37.0\%\\
all & 0.400 & 0 & 0.389 & 33.4\% & 0.282 & 48.8\%\\
\bottomrule
\end{tabular}
\end{center}

Capture is the provider's mark-out per request, filled or not, in basis points (the mechanism in
\cref{fig:hf:fx-market-making:lastlook}). Last look pays only on the toxic tier: it rejects every informed
request there, and 19\% of the uninformed ones with them. On the other tiers it gives up more in rejected good trades than it saves, and the
symmetric check, which rejects moves in the provider's favour too, gives up more still. What last look costs the clients (rejections and the
information in them) is the subject of the FX Global Code's principle on it; the arithmetic here is the provider's side.

\section{Internalise or hedge}

Each accepted trade adds to the provider's inventory. It can hedge at once on a primary venue, paying that venue's spread, or internalise: keep the
position and let later client trades offset it, helped by skewing its stream so that clients are more likely to take the side that reduces it. The
chapter's day (\cref{lst:hf:fx-market-making:internalise}): 20,000 client trades of \$1 million, 0.4 basis point earned on each, hedging at 0.25
basis point, the rate moving 60 basis points over the day, and a charge on the inventory P\&L's variance. The skew makes the inventory
mean-revert; its half-life, in client trades, decides the answer.

\begin{omfigure}
\begin{tikzpicture}
\begin{axis}[width=11cm,height=5.4cm,axis y line*=right,axis x line=none,xmode=log,log basis x=10,xmin=50,xmax=2500,ymin=-250,ymax=800,
  ylabel={\small objective (\$ thousand a day)},tick label style={font=\footnotesize},table/col sep=comma]
\addplot[omDef,thick,dashed,mark=square*,mark size=1.3pt] table[x=halflife,y=internalise]{figdata/hft/21-fx-market-making/hedge.csv};
\label{plot:hf:fxint}
\addplot[gray,thin] table[x=halflife,y=hedge]{figdata/hft/21-fx-market-making/hedge.csv};
\label{plot:hf:fxhedge}
\end{axis}
\begin{axis}[width=11cm,height=5.4cm,axis y line*=left,xmode=log,log basis x=10,xmin=50,xmax=2500,xtick={50,100,200,500,1000,2000},xticklabels={50,100,200,500,1000,2000},xlabel={\small inventory half-life (client
  trades)},ylabel={\small best share hedged at once (\%)},ymin=0,ymax=105,tick label style={font=\footnotesize},grid=major,grid style={gray!25},
  table/col sep=comma,legend style={font=\scriptsize,at={(0.5,-0.3)},anchor=north},legend columns=3]
\addplot[omThm,thick,mark=*,mark size=1.3pt] table[x=halflife,y=best]{figdata/hft/21-fx-market-making/hedge.csv};
\addlegendentry{best hedge share (left)}
\addlegendimage{/pgfplots/refstyle=plot:hf:fxint}\addlegendentry{internalise all (right)}
\addlegendimage{/pgfplots/refstyle=plot:hf:fxhedge}\addlegendentry{hedge all (right)}
\end{axis}
\end{tikzpicture}

\omcaption{The share of each client trade best hedged at once on the primary venue (left), and the day's objective (spread less hedging costs less
a charge of 0.002 per squared hundred dollars of inventory P\&L variance, right) when internalising everything or hedging everything, against the
inventory's half-life under the stream's skew (log scale); 20,000 trades of \$1 million, 0.4 basis point earned, 0.25 paid to hedge. Data:
\texttt{hf\_fx.hedging}.}\label{fig:hf:fx-market-making:hedge}
\end{omfigure}

Without skew the inventory is a random walk and the best policy hedges 95\% of every trade. With a half-life of about 1,150 trades or more it still
hedges 95\%; at 870, 80\%; at 770, half; at 690 or less, almost nothing (\cref{fig:hf:fx-market-making:hedge}). Once the flow offsets itself fast
enough, internalising earns far more than hedging: \$588,000 a day at a half-life of 350 trades against \$300,000 for hedging everything. The switch
is abrupt for a reason worth knowing. Without skew the inventory's variance grows with the square of the unhedged share $1-h$, and the objective
$S-hC-\lambda(1-h)^2V$ has its maximum at $h^\ast=1-C/(2\lambda V)$: 97\% here. With the skew, a larger inventory is pulled back harder, and the
variance grows only about in proportion to $1-h$: the objective is then nearly linear in $h$, and its maximum sits at a corner.

\section{Crosses and emerging markets}

Crosses and emerging-market currencies are where internalisation and synthesis matter most: their direct markets are thin, their legs through the
dollar are liquid, and a provider with flow in many pairs can offset one client's euro-yen against another's euro-dollar and a third's dollar-yen.
Non-deliverable forwards (One Quant Book 2, chapter 18) extend electronic making to currencies with capital controls; they settle in dollars on a
fixing, so the provider's risk runs to the fixing date and its hedge is another non-deliverable forward or the onshore market where it has access.

\section{Strategy files}

\begin{strategyfile}{Primary-venue FX market making}\label{strat:hf:fx-market-making:primary}
\sfield{Who pays you, and why} Banks and funds that take liquidity on the interdealer venues.
\sfield{Instruments and venues} Major pairs on the primary interdealer venues.
\sfield{Signal} A \omterm{def:hf:fair-value:consolidated}{consolidated fair price} from all venues, weighted by noise and staleness.
\sfield{Sizing and execution} Quote at the touch with inventory skew; firm quotes, no last look.
\sfield{Costs} Venue fees; adverse selection from faster takers.
\sfield{How it dies} Speed: algorithmic takers remove the arbitrage opportunities and impose adverse selection on slower providers.
\sfield{Horizon, capacity, infrastructure} Microseconds to seconds; co-location at each venue.
\sfield{Backtest honestly} Each venue's data at its own delay; queue positions.
\sfield{Sources} Chaboud, Chiquoine, Hjalmarsson and Vega (2014).
\end{strategyfile}

\begin{strategyfile}{Disclosed streaming to aggregators with last look}\label{strat:hf:fx-market-making:stream}
\sfield{Who pays you, and why} Clients who value a streamed price in their own aggregator, by tier.
\sfield{Instruments and venues} Streams to banks, brokers, platforms and corporates.
\sfield{Signal} The \omterm{def:hf:fair-value:fair}{fair price}; each tier's measured toxicity.
\sfield{Sizing and execution} Price each tier to its mark-outs; apply last look only where it pays (the toxic tier in the model).
\sfield{Costs} Rejects that clients remember; the hedge.
\sfield{How it dies} Clients who move flow to firm liquidity; the Code's principles on last look and its disclosure.
\sfield{Horizon, capacity, infrastructure} Milliseconds; hundreds of streams.
\sfield{Backtest honestly} Rejections as the clients saw them, and the flow lost because of them.
\sfield{Sources} One Quant Book 2, chapter 15; this chapter.
\end{strategyfile}

\begin{strategyfile}{Cross-rate synthesis quoting}\label{strat:hf:fx-market-making:cross}
\sfield{Who pays you, and why} Clients trading crosses whose direct markets are thin.
\sfield{Instruments and venues} A cross and its two legs through a common currency.
\sfield{Signal} The \omterm{def:hf:fx-market-making:cross}{synthetic cross} from the legs' bids and offers.
\sfield{Sizing and execution} Quote the cross inside its synthetic width when flow can be internalised; hedge in the legs or the direct market,
whichever is cheaper.
\sfield{Costs} Two legs' spreads when hedging (2.6 pips on the example).
\sfield{How it dies} One-sided cross flow that cannot be internalised.
\sfield{Horizon, capacity, infrastructure} Seconds; data on both legs and the direct cross.
\sfield{Backtest honestly} Leg prices executable at the cross trade's time.
\sfield{Sources} This chapter.
\end{strategyfile}

\begin{strategyfile}{Non-deliverable forward electronic making}\label{strat:hf:fx-market-making:ndf}
\sfield{Who pays you, and why} Investors and companies hedging currencies they cannot deliver.
\sfield{Instruments and venues} Non-deliverable forwards on electronic platforms.
\sfield{Signal} Onshore rates where observable, the fixing's history, related liquid currencies.
\sfield{Sizing and execution} Quote tiered streams; internalise across tenors and clients; hedge the residual in other forwards.
\sfield{Costs} Wide hedging markets; fixing risk.
\sfield{How it dies} Capital-control changes and fixing disputes.
\sfield{Horizon, capacity, infrastructure} Days to months to the fixing.
\sfield{Backtest honestly} The fixing as published, with its source's history.
\sfield{Sources} One Quant Book 2, chapter 18.
\end{strategyfile}

\section{Tutorial: a hundred prices at once}\label{tut:hf:fx-market-making:tut}

\begin{tutorial}
\textbf{Goal.} Build a \omterm{def:hf:fair-value:consolidated}{consolidated fair price}, stream it to three tiers with and without last look, and decide how much of the flow to internalise.
\textbf{End state:} the table and \cref{fig:hf:fx-market-making:hedge}.
\begin{tutsteps}
\item \textbf{The \omterm{def:hf:fair-value:fair}{fair price}}: inverse-variance weights with staleness.
\omcode{code/firm/fxlp/firm_fxlp.py}{34}{38}{Each venue's noise plus the rate's variance since its last update.}\label{lst:hf:fx-market-making:fair}
\item \textbf{The streams} on Book 2's \texttt{firm.lastlook}.
\omcode{code/firm/fxlp/firm_fxlp.py}{57}{71}{Each tier's requests, the last-look check and its transaction-cost analysis, weighted by the tier's share.}\label{lst:hf:fx-market-making:stream}
\item \textbf{Internalise or hedge} with a skew that brings offsetting flow.
\omcode{code/firm/fxlp/firm_fxlp.py}{81}{98}{The inventory under a hedge share and a skew; spread, hedge cost and the risk charge.}\label{lst:hf:fx-market-making:internalise}
\item \textbf{Sweep} the hedge share for each skew (\texttt{hf\_fx.hedging}).
\end{tutsteps}
\textbf{What to change next.} Let informed clients cluster in time; add a second pair and internalise across them; run the streams on
\texttt{firm.exchsim}'s venues with their own latencies.
\end{tutorial}

\section{Build: the FX liquidity provider}\label{bld:hf:fx-market-making:fxlp}

\begin{build}
\bfield{Purpose} Price FX from many venues, stream tiered prices with or without last look, and decide internalisation.
\bfield{Interface} \texttt{fair\_price(mids, noise\_bp, ages\_ms, sigma)}, \texttt{venue\_errors}, \texttt{stream(tiers, hold, threshold, policy,
sigma, n, seed)}, \texttt{internalise(h, seed, trades, size, sigma\_day, half\_bp, hedge\_bp, lam, skew)}, \texttt{best\_hedge\_share},
\texttt{synthetic\_cross}. Built on \texttt{firm.lastlook} (Book 2).
\bfield{Rules} Basis points of the rate; milliseconds; amounts in basis points times millions.
\bfield{Acceptance tests} \texttt{code/firm/fxlp/tests/}: weights by hand and a fresh, precise venue dominating; the consolidated price beats the
primary; no rejects without last look, all informed rejected with the asymmetric check; hedging everything leaves no inventory risk; a stronger skew
lowers the inventory; the \omterm{def:hf:fx-market-making:cross}{synthetic cross} by hand, direct and inverted.
\bfield{Stretch} Clustered informed flow; several pairs internalised together; streams on \texttt{firm.exchsim}.
\end{build}

\begin{omsources}
\item A.~P.~Chaboud, B.~Chiquoine, E.~Hjalmarsson, C.~Vega, Rise of the machines: algorithmic trading in the foreign exchange market, \emph{Journal
of Finance} 69(5), 2014, 2045--2084.
\item Bank for International Settlements, Triennial Central Bank Survey, preliminary results, press release, 30 September 2025.
\end{omsources}

\section{Exercises}

\begin{exercise}[$\star$]\label{exo:hf:fx-market-making:1}
Two venues show 1.08500 and 1.08503 with noise of 0.3 and 0.6 basis point, both fresh. What is the \omterm{def:hf:fair-value:fair}{fair price}?
\end{exercise}

\begin{exercise}[$\star$]\label{exo:hf:fx-market-making:2}
Build euro-yen from euro-dollar 1.0850/1.0851 and dollar-yen 150.20/150.21. How wide is it in pips?
\end{exercise}

\begin{exercise}[$\star$]\label{exo:hf:fx-market-making:3}
At 0.4 basis point on \$20 billion a day, what does the stream earn before hedging? What does hedging everything at 0.25 cost?
\end{exercise}

\begin{exercise}[$\star\star$]\label{exo:hf:fx-market-making:4}
Why does last look raise the toxic tier's capture and lower the others'?
\end{exercise}

\begin{exercise}[$\star\star$]\label{exo:hf:fx-market-making:5}
Why is the best hedge share almost always either near 0 or near 95\%?
\end{exercise}

\begin{exercise}[$\star\star$]\label{exo:hf:fx-market-making:6}
What did Chaboud and co-authors find about who removed triangular arbitrage?
\end{exercise}

\begin{exercise}[$\star\star\star$]\label{exo:hf:fx-market-making:7}
\emph{Coding.} Rerun \texttt{firm.fxlp.stream} with a hold of 20 milliseconds instead of 5. What happens to the rejects of uninformed tier-3 clients?
\end{exercise}

\begin{exercise}[$\star\star\star$]\label{exo:hf:fx-market-making:8}
\emph{Find the flaw.} ``Our internalisation rate is 90\%, so our hedging costs are a tenth of our competitors'; we are the cheapest provider.''
\end{exercise}

\section{Problem: A Hundred Prices at Once}

\begin{problem}[{Weekend problem --- a hundred prices at once}]\label{pb:hf:fx-market-making:1}
A non-bank liquidity provider prices euro-dollar from three venues, streams it to three tiers, and decides what to hedge.

\medskip\noindent\textbf{Part I --- The market.}
\begin{enumerate}
\item Summarise the dated box.
\item What did Chaboud and co-authors find about algorithmic trading in FX?
\item How is the \omterm{def:hf:fair-value:fair}{fair price} built, and how much better is it than the primary?
\item Define a \omterm{def:hf:fx-market-making:cross}{synthetic cross} and give the example's width.
\end{enumerate}

\medskip\noindent\textbf{Part II --- Streams.}
\begin{enumerate}[resume]
\item Describe the tiers and the last-look check.
\item Give capture and rejects without last look and with the asymmetric check, per tier.
\item Why is the symmetric check worse for the provider?
\item What does last look cost the clients?
\end{enumerate}

\medskip\noindent\textbf{Part III --- Inventory.}
\begin{enumerate}[resume]
\item Describe the internalisation model and the skew.
\item Give the best hedge share at half-lives of 1,150, 870, 770 and 690 trades.
\item What does internalising earn at a half-life of 350 trades?
\item Why is the switch abrupt?
\end{enumerate}

\medskip\noindent\textbf{Part IV --- The verdict.}
\begin{enumerate}[resume]
\item State the \emph{named result}: the liquidity provider's capture and reject rate with and without last look, and the hedge share at which
internalising stops paying.
\item Which tiers would you give last look, and which firm prices?
\item How would you measure the inventory's half-life in real flow?
\item What makes crosses and emerging markets different?
\item Which strategy file is most exposed to regulation?
\item What does a provider with flow in many pairs gain?
\item How would clustered informed flow change the answers?
\item In one sentence: what does an FX liquidity provider sell?
\end{enumerate}
\end{problem}

\section{Interview questions}

\begin{interviewq}[$\star$ \iqroles{trader}]\label{iq:hf:fx-market-making:1}
A client lifts your euro-dollar offer for \$50 million. Hedge now or wait? What do you look at?
\end{interviewq}

\begin{interviewq}[$\star\star$ \iqroles{researcher}]\label{iq:hf:fx-market-making:2}
How would you measure each client's toxicity, and what would you do with the measure?
\end{interviewq}

\begin{interviewq}[$\star\star$ \iqroles{developer}]\label{iq:hf:fx-market-making:3}
Design a system that streams 200 tiered prices on 30 pairs and updates them within 50 microseconds of a venue tick.
\end{interviewq}

\begin{interviewq}[$\star\star$ \iqroles{risk}]\label{iq:hf:fx-market-making:4}
Your internalisation rate rose from 70\% to 95\% in a month. What could explain it, and what could go wrong?
\end{interviewq}

\begin{interviewq}[$\star\star$ \iqroles{trader}]\label{iq:hf:fx-market-making:5}
Explain why a symmetric last-look check is fairer to clients than an asymmetric one, and what it costs the provider.
\end{interviewq}

\begin{interviewq}[$\star\star\star$ \iqroles{researcher}]\label{iq:hf:fx-market-making:6}
With spread $S$, hedging cost $hC$ and a risk charge $\lambda\,v(h)$, find the optimal hedge share when $v(h)=(1-h)^2V$ and when
$v(h)=(1-h)V'$.
\end{interviewq}
